\section{Related Work}
\label{sec:related}

\subsection{RLVR and On-Policy Distillation} 
Reinforcement learning with verifiable rewards (RLVR)~\citep{shao2024deepseekmath,yu2026dapo,liu2025understanding,zheng2025group,pan2025d} has emerged in recent years as a central post-training paradigm for large language models, achieving substantial gains on reasoning-intensive benchmarks. Representative algorithms such as GRPO~\citep{shao2024deepseekmath} dispense with the critic network of classical actor-critic methods and instead estimate advantages by normalizing rewards across a group of rollouts sampled from the same query. The reward itself is supplied by an external verifier and is therefore noise-free on the tasks it covers. The granularity of supervision, however, is at the level of the entire response: a single scalar advantage is broadcast uniformly to every token in the rollout, providing no fine-grained credit assignment to the intermediate reasoning steps that determine correctness. 

On-policy distillation (OPD)~\citep{gu2024minillm,agarwal2024policy} has emerged in response to this limitation. By computing a per-token distribution gap between a teacher and the student along the student's own sampled trajectories, OPD yields a dense, token-level learning signal that complements the sparse trajectory-level reward of RLVR. The dominant choice for measuring this gap is the reverse-KL divergence between teacher and student, giving rise to the canonical OPD objective:
\begin{equation}
\mathcal{L}_{\text{OPD}}(\theta) = \mathbb{E}_{y \sim \pi_S} \left[ \frac{1}{|y|} \sum_{t=1}^{|y|} \mathrm{KL}\left( \pi_S(\cdot \mid x, y_{<t}) \,\|\, \mathrm{sg}[\pi_T(\cdot \mid x, y_{<t})] \right) \right],
\end{equation}
where $\mathrm{sg}[\cdot]$ denotes the stop-gradient operator and $|y|$ is the length of the rollout in tokens. It has been shown that this objective is equivalent, in the policy-gradient sense, to REINFORCE~\citep{williams1992simple} with per-token reward $r_t=\log \pi_T(y_t\mid x,y_{<t}) - \log \pi_S(y_t\mid x,y_{<t})$ ~\citep{gu2024minillm,thinkingmachines2025opd}. 

A line of work has explored alternative divergence choices to address the trade-off between mode coverage and mode seeking inherent in reverse-KL: forward-KL offers stronger mode coverage at the cost of mass-spreading, while Jensen-Shannon Divergence~\citep{agarwal2024policy} and skew KL~\citep{ko2024distillm} formulations interpolate between the two extremes. Other extensions consider OPD under black-box access to the teacher~\citep{ye2025black}, and propose training-efficient variants that reduce the per-step distillation cost~\citep{wu2026lightning}. At industrial scale, MiMo-V2-Flash~\citep{xiao2026mimo} integrates the dense token-level advantage $\log \pi_T(y_t\mid x,y_{<t}) - \log \pi_S(y_t\mid x,y_{<t})$ from OPD with the outcome-based verifier reward $A_\mathrm{ORM}$ of RLVR in its post-training pipeline. More recent investigations~\citep{li2026rethinking,fu2026revisiting,zhu2026many} have shifted attention from objective design to the conditions under which OPD itself succeeds: the teacher and student should share compatible reasoning styles, the teacher must possess capabilities genuinely absent from the student, response length must fall within a productive regime, and approximations such as top-$k$ truncation should preserve unbiasedness of the gradient estimator.

\subsection{On-Policy Self-Distillation}
Despite the rapid progress of on-policy distillation, the paradigm faces a series of practical obstacles. First, most OPD methods require white-box access to the teacher's token-level logits, which precludes the use of strong closed-source models as teachers. Second, the teacher and student must share an identical vocabulary, a condition rarely satisfied across model families and one that restricts OPD in practice to the distillation of smaller models from larger counterparts within the same series. Third, serving a separate teacher model throughout RL training imposes substantial memory and latency overhead. To address these limitations, on-policy self-distillation (OPSD) has emerged~\citep{zhao2026self,han2026adaptive,yu2026preference,hao2026self,heakl2026cepo}. In OPSD, the teacher and the student are the same model; the teacher is simply granted access to privileged context unavailable to the student, such as a ground-truth solution or terminal feedback from the environment. The standard OPSD objective takes the following dense distillation form:
\begin{equation}
\mathcal{L}_{\text{OPSD}}(\theta) = \mathbb{E}_{y \sim \pi_S} \left[ \frac{1}{|y|} \sum_{t=1}^{|y|} \mathrm{KL}\left( \pi_S(\cdot \mid x, y_{<t}) \,\|\, \mathrm{sg}[\pi_T(\cdot \mid x, y_c^*, y_{<t})] \right) \right],
\end{equation}
where $T$ and $S$ denotes the same model, and $y_c^*$ is the privileged context. Inherited from OPD, the OPSD objective is equivalent to reinforcement learning with per-token advantage $A_{t} = \log \pi_T(y_t \mid x, y_c^*, y_{<t}) - \log \pi_S(y_t \mid x, y_{<t})$. 

Recent work has explored different instantiations of the OPSD framework, most of which adopt \emph{dense} per-token distillation. OPSD~\citep{zhao2026self} performs dense distillation using the canonical forward-KL form. While it observes the dominance of stylistic tokens in the per-token signal, it addresses this only through an ad-hoc per-token pointwise KL-clipping mechanism, which fails to resolve the issue at its root and leaves training unstable. SDPO~\citep{hubotter2026reinforcement} likewise uses dense distillation, but replaces forward KL with the Jensen--Shannon divergence to balance mode-seeking and mode-covering behavior in the per-token gap. SRPO~\citep{li2026unifying} observes that the privileged context contributes little additional signal on already-correct rollouts, and proposes a sample-level routing mechanism in which correct samples are optimized via GRPO while incorrect ones are optimized via SDPO's dense objective. In contrast to all of the above, RLSD~\citep{yang2026self} departs from dense distillation entirely: rather than matching full token-level distributions, it relies on a \emph{sampled-token} signal (which RLCSD also adopts) and argues that this OPSD-like signal should serve only as a modulation of $A_{\text{ORM}}$ rather than a replacement. Accordingly, it designs an algorithm in which the verifier-derived $A_{\text{ORM}}$ determines the direction of every update while the sampled-token signal modulates only its magnitude.

However, all of these methods rely exclusively on \emph{one-sided positive privileged contexts} (e.g., dataset ground-truth solutions or correct in-group rollouts) when constructing the teacher distribution, and therefore inherit the privilege-induced style drift identified earlier as a systematic confound in their token-level signal. \textbf{To address this limitation, we propose RLCSD, which removes the drift by construction through a symmetric contrastive cancellation between correct and incorrect privileged contexts}.



